\begingroup
\emergencystretch=3em
El desarrollo del Trabajo Terminal 1 permitió delimitar una propuesta para apoyar la selección de alineaciones de la WNBA mediante el análisis de coincidencias individuales y el rendimiento colectivo. El resultado de esta etapa es una especificación metodológica y de datos que distingue qué información puede conocerse antes de una consulta, qué resultados se pretende estimar y cómo se compararán las alternativas. No constituye todavía una comprobación de precisión predictiva ni evidencia de mejora en decisiones reales de sustitución; esas conclusiones dependerán de la implementación y evaluación previstas para el Trabajo Terminal 2.

Uno de los aprendizajes centrales fue que la disponibilidad de estadísticas no equivale a disponer directamente de evidencia para todas las preguntas deportivas. El registro de eventos puede permitir reconstruir quiénes coincidieron en cancha, pero esa coincidencia no identifica necesariamente quién defendió a quién. Por ello, el análisis denominado individual se delimitó como una estimación de producción estadística durante la presencia de una rival, sin atribuirle el efecto causal de un marcaje. Esta distinción reduce el alcance de las afirmaciones posibles, pero también establece una base verificable para trabajar con los datos disponibles.

La ausencia de seguimiento espacial y de asignaciones defensivas no se resolverá únicamente mediante nuevos algoritmos. El procesamiento podrá recuperar intervalos, posesiones y acciones registradas, mientras que la información no observada seguirá siendo una limitación. En consecuencia, no se prometerá identificar a la mejor defensora directa ni medir todas las contribuciones sin balón. La alternativa adoptada consiste en formular preguntas compatibles con la evidencia: estimar producción individual condicionada a coincidencias y valorar el balance colectivo esperado de quintetos frente a un rival definido.

La revisión de los trabajos permitió situar la propuesta dentro de una línea que ya estudia perfiles estadísticos, interacciones entre integrantes y desempeño de alineaciones \cite{yamada2024offensivelineupanalysisbasketball,josephs2025hypergraphadjustedplusminus}. Por tanto, no se sostiene que el análisis de quintetos o el uso de perfiles sean aportaciones inéditas por sí mismos. La diferencia que se propone evaluar es su articulación en dos fases: generar puntuaciones individuales para parejas comparables y comprobar si enriquecen la predicción colectiva utilizada para seleccionar una alineación. La aportación dependerá de demostrar utilidad respecto a referencias sencillas, no de combinar técnicas o aumentar la cantidad de variables.

También se precisó que predecir y optimizar son tareas distintas. El predictor estimará el Net Rating de cada candidato; la selección buscará el mayor valor predicho entre quintetos que respeten disponibilidad, composición y límites de sustituciones. Una buena posición individual no obligará a incluir a una jugadora si otra combinación presenta mejor valoración colectiva. A su vez, encontrar el máximo entre las predicciones no demostrará que se haya encontrado el mejor quinteto real, particularmente cuando existan pocas posesiones observadas o combinaciones sin historial.

La separación entre tablas maestras, antecedentes y entradas predictivas constituye otra decisión relevante del diseño. Los resultados del partido objetivo se reservarán para construir las respuestas y evaluar errores; no se rellenarán como si fueran conocidos al recomendar. Los datos faltantes requerirán distinguir ausencia de historial, incompatibilidad de roles e información incompleta. Estas precauciones, junto con los cortes cronológicos y las predicciones temporales fuera de muestra, permitirán evaluar el procedimiento sin utilizar resultados futuros en su preparación \cite{sklearnCommonPitfallsPonderacion,cawley2010selection}.

Para el Trabajo Terminal 2, cuyo periodo previsto de desarrollo es de aproximadamente seis meses, se priorizará una implementación completa y acotada del procedimiento estático. Primero se verificará la cobertura de las fuentes y se construirán los intervalos, las tablas maestras y los antecedentes reproducibles. Después se compararán las respuestas individuales previstas, se seleccionará la ponderación temporal dentro del desarrollo y se implementarán los modelos iniciales de ambas fases. La evaluación colectiva incluirá la comparación con y sin puntuaciones de fase 1 y la selección de quintetos factibles mediante enumeración. Finalmente, el demostrador local permitirá introducir escenarios, ejecutar consultas y registrar resultados y pruebas funcionales.

El clustering se tratará como un enriquecimiento cuya continuidad dependerá de su aportación, no como un resultado que deba conservarse a cualquier costo. Los algoritmos adicionales y las características derivadas de hipergrafos quedarán condicionados a cobertura, resultados de validación y tiempo disponible. No se comprometerá la implementación de todas las alternativas para considerar concluido el trabajo. Si alguna configuración no mejora las referencias o su evidencia resulta insuficiente, se documentará ese resultado y se conservará el procedimiento más sencillo respaldado por los experimentos. Tampoco se reutilizará la prueba final para presentar una nueva selección como evaluación independiente.

La relevancia del proyecto para el baloncesto femenil se plantea en términos concretos: organizar evidencia de la WNBA, hacer explícitos los criterios de comparación y evaluar una herramienta reproducible de apoyo analítico. El procedimiento podría facilitar la exploración de alternativas y sus antecedentes, sin sustituir el conocimiento técnico del cuerpo entrenador. Su valor deberá sustentarse en precisión, cobertura, trazabilidad y límites informados; no se afirmará un impacto deportivo general a partir de una demostración retrospectiva.

Como líneas posteriores al alcance comprometido podrán considerarse datos de asignación defensiva, seguimiento espacial, actualización durante el partido y evaluación con personal especializado. Estas extensiones requerirían nuevas fuentes, recursos y protocolos, por lo que no serán condiciones de entrega del Trabajo Terminal 2. La conclusión esperada de esa etapa será establecer, mediante resultados verificables, en qué condiciones la integración propuesta aporta información útil para comparar alineaciones y en cuáles la incertidumbre o la falta de datos impiden sostener una recomendación fiable.
\endgroup
